\documentclass[a4paper, 10pt]{article}
\usepackage{scrextend}
\changefontsizes{9pt}
% This file lives at vendor/template/templates/style-guide/ once vendored
% into a subject repo -- course-config.tex sits at the subject repo root (4
% levels up), while CEDT-Assignment-style.sty sits at the vendor/template/
% root itself (2 levels up; no need to spell "vendor/template" again, we're
% already inside it).
\input{../../../../course-config}
\usepackage{../../CEDT-Assignment-style}

\begin{document}
\subject
\asgntitle{}{CEDT-Assignment-style Reference}{Every macro and environment this style provides, with a worked example}{\StudentID\ \StudentName}{}
% \setfooter overrides the footer \asgntitle just set above -- useful when you
% want a shorter running footer than the full printed title.
\setfooter{Style Reference}

\tableofcontents

% ================================================================================ %
\section{Math shortcuts}
% ================================================================================ %
Common number sets and delimiters get short macros so you don't have to
remember the \texttt{amsmath} spelling every time:
\begin{itemize}
	\item Sets: $\R$, $\C$, $\Z$, $\N$ -- \verb|\R \C \Z \N|
	\item Delimiters: $\abs{x}$, $\norm{v}$, $\paren{a+b}$, $\sqbracket{a,b}$ --
	      \verb|\abs \norm \paren \sqbracket|
\end{itemize}
For example: $\forall x \in \R,\ \abs{x} \geq 0$, and
$\norm{v} = \sqrt{\paren{v_1^2 + v_2^2}}$.

% ================================================================================ %
\section{Problems and sub-problems}
% ================================================================================ %
\begin{problem}
\verb|problem| auto-numbers itself. Sub-items nest under \verb|subproblems|
(decimal, e.g. 2.1) by default:
\begin{subproblems}
	\item Decimal sub-items look like this.
	\item A sub-item can go one level deeper with \verb|Tsubsubproblems|:
	      \begin{Tsubsubproblems}
		      \item A "T"-prefixed third nesting level -- adjust the fixed "T" in
		            the style file if it doesn't fit your subject.
		      \item Another one at the same depth.
	      \end{Tsubsubproblems}
\end{subproblems}
Or letter-labelled with \verb|subproblems_alpha| instead:
\begin{subproblems_alpha}
	\item Like this.
	\item And this.
\end{subproblems_alpha}
\end{problem}

\begin{problem}[10]
Passing a number, e.g. \verb|\begin{problem}[10]|, sets/reuses that exact
problem number instead of auto-incrementing -- handy when your numbering has
to match a printed assignment sheet. Auto-numbering after this continues from
whatever the highest-used number is, skipping any manually reused ones.
\end{problem}

% ================================================================================ %
\section{Theorem-style environments}
% ================================================================================ %
Standard numbered \texttt{amsthm} environments are set up too --
\texttt{theorem}/\texttt{lemma}/\texttt{proposition}/\texttt{corollary} share
one counter, \texttt{definition} numbers separately, and \texttt{remark} is
unnumbered:

\begin{definition}
	An integer $n$ is \emph{even} if $n = 2k$ for some $k \in \Z$.
\end{definition}

\begin{theorem}[Fundamental theorem of arithmetic]
	Every integer greater than $1$ has a unique prime factorization, up to the
	order of factors.
\end{theorem}

\begin{lemma}
	If $n^2$ is even, then $n$ is even.
\end{lemma}

\begin{proposition}
	The sum of two even integers is even.
\end{proposition}

\begin{corollary}
	If $n$ is even, then $n^2$ is even.
\end{corollary}

\begin{remark}
	These come straight from \texttt{amsthm} -- nothing custom beyond the
	numbering/sharing scheme set up in the style file.
\end{remark}

% ================================================================================ %
\section{Solution, proof, and to-submit boxes}
% ================================================================================ %
\begin{problem}
Worked-through answers that aren't the graded deliverable pair with
\verb|solution|.
\end{problem}
\begin{solution}
	The full derivation goes here.

	A second paragraph indents correctly too -- \texttt{CEDT-Assignment-style.sty}
	restores \verb|\parindent| inside every breakable box below (tcolorbox's
	\texttt{breakable} library, needed so long boxes can split across a page,
	zeroes it internally otherwise).
\end{solution}

\begin{problem}
A problem that must be proven pairs with \verb|proofbox|.
\end{problem}
\begin{proofbox}
	The proof goes here.

	A second paragraph, indented the same way. \hfill\qedsymbol
\end{proofbox}

\begin{problem}
A problem whose final answer must be called out for grading pairs with
\verb|tosubmit| -- its corner tab reads "TO SUBMIT | Solution" rather than
just "TO SUBMIT".
\end{problem}
\begin{tosubmit}
	This is what a grader should look at directly.

	A second paragraph here indents too.
\end{tosubmit}

% ================================================================================ %
\section{Code blocks}
% ================================================================================ %
\verb|codingbox| defaults to Python, but any language \texttt{listings} ships
a \verb|\lstdefinelanguage| for works via its \texttt{language} option:
\begin{codingbox}
    print("Python is the default language")
\end{codingbox}
\begin{codingbox}[language=C++]
#include <iostream>

int main() {
    // Compute the first 10 Fibonacci numbers
    int a = 0, b = 1;
    for (int i = 0; i < 10; ++i) {
        std::cout << a << " ";
        int next = a + b;
        a = b;
        b = next;
    }
    return 0;
}
\end{codingbox}
\begin{codingbox}[language=bash]
#!/usr/bin/env bash
# Print the first 10 Fibonacci numbers
a=0
b=1
for _ in $(seq 1 10); do
    echo -n "$a "
    next=$((a + b))
    a=$b
    b=$next
done
\end{codingbox}

% ================================================================================ %
\section{Figures}
% ================================================================================ %
\verb|\fig[width]{path}{caption}{label}| centers an image in a numbered,
captioned, labelled \texttt{figure} float (width defaults to
\verb|0.7\linewidth|):
\fig{../template/images/fibonacci-example.png}{A numbered, captioned figure.}{fib-numbered}

The starred form, \verb|\fig*[width]{path}{}{}|, drops the caption/label
entirely -- just the centered image (the last two arguments are still
required syntactically, but ignored, so pass \verb|{}{}|):
\fig*[0.4\linewidth]{../template/images/fibonacci-example.png}{}{}

% ================================================================================ %
\section{Tables}
% ================================================================================ %
\verb|\begin{simpletable}{<colspec>}{<caption>}{<label>}| wraps a
\texttt{booktabs}-style table with a numbered caption/label -- write rows as
normal and add \verb|\midrule| after the header yourself:
\begin{simpletable}{lcc}{Truth table for logical AND}{and-truth}
	$A$ & $B$ & $A \wedge B$ \\
	\midrule
	0   & 0   & 0             \\
	0   & 1   & 0             \\
	1   & 0   & 0             \\
	1   & 1   & 1             \\
\end{simpletable}

% ================================================================================ %
\section{Multiple choice}
% ================================================================================ %
\begin{problem}
\verb|\choices[<num choices>]{choice A, choice B, ...}| lays out options,
auto-picking 4, 2, or 1 columns depending on how long the choices are (never
more columns than there are choices):
\begin{subproblems}
	\item Short choices sit four to a row.
	      \choices{True, False, Maybe, Unknown}
	\item Longer choices wrap down to two columns automatically.
	      \choices{The probability increases, The probability decreases, The probability stays the same, Not enough information to tell}
	\item Pass the choice count explicitly when it isn't 4.
	      \choices[6]{A, B, C, D, E, F}
\end{subproblems}
\end{problem}

% ================================================================================ %
\section{Spacing helper}
% ================================================================================ %
\verb|\myspace[h|\emph{or}\verb| v][amount][unit]| wraps \verb|\hspace|/\verb|\vspace|;
all three arguments are optional and default to \verb|\myspace| = \verb|\vspace{4mm}|.
For example, \verb|A\myspace[h][1][cm]B| gives A\myspace[h][1][cm]B (a 1cm gap
between A and B on the same line).

% ================================================================================ %
%                                     Appendix                                     %
% ================================================================================ %
% \myappendix{path/to/code.pdf}{Title} appends a title page + the full PDF
% (e.g. a notebook exported with `make nb-pdf`). Safe to point at a real file
% here since this reference doc, unlike assignment-template.tex, never gets
% copied into a fresh assignment folder. (../template/code/ -- the demo
% notebook lives next to assignment-template.tex, not next to this file.)
\myappendix{../template/code/notebook-template.pdf}{notebook-template | Code Solution}

\end{document}
